\sectionheading{5.5 মধ্যম স্তরের ভাষা (Middle Level Language)}

Programming Language-এর কিছু ভাষায় High Level Language-এর
সহজে Program লেখার সুবিধা এবং Low Level Language-এর
Hardware-এর কাছাকাছি কাজ করার ক্ষমতা—উভয় বৈশিষ্ট্যই পাওয়া যায়।
এ ধরনের ভাষাকে \textbf{Middle Level Language} বলা হয়।

%================================================
% Definition
%================================================
\booksubheading{মধ্যম স্তরের ভাষা কাকে বলে?}

\begin{definitionbox}

\textbf{\color{darkTeal}সংজ্ঞা}

যে Programming Language-এ \textbf{High Level Language ও Low Level
Language—উভয়ের গুরুত্বপূর্ণ বৈশিষ্ট্য বিদ্যমান থাকে}, তাকে
\textbf{মধ্যম স্তরের ভাষা (Middle Level Language)} বলে।

\end{definitionbox}

\textbf{সহজভাবে মনে রাখবে—}

\begin{center}

\textbf{
High Level-এর সুবিধা
\(\;+\;\)
Low Level-এর নিয়ন্ত্রণ
\(\;=\;\)
Middle Level Language
}

\end{center}

এখানে ``Middle'' বলতে ভাষার মান মাঝারি বোঝায় না; বরং High Level
ও Low Level—দুই ধরনের বৈশিষ্ট্যের সমন্বয় বোঝায়।

%================================================
% Examples
%================================================
\booksubheading{Middle Level Language-এর উদাহরণ}

Middle Level বৈশিষ্ট্যসম্পন্ন ভাষার উদাহরণ হিসেবে
\textbf{C, C++} ইত্যাদির নাম উল্লেখ করা হয়।

HSC পর্যায়ে \textbf{C Language}-কে Middle Level Language-এর
সবচেয়ে পরিচিত উদাহরণ হিসেবে আলোচনা করা হয়।

\begin{examplebox}{সহজ উদাহরণ}

C ভাষা দিয়ে যেমন সাধারণ হিসাবের Program লেখা যায়, তেমনি
Pointer, Memory Address ও Bitwise Operator ব্যবহার করে
Hardware-এর কাছাকাছি কাজও করা যায়।

অর্থাৎ—

\begin{center}

\textbf{
General Programming
\(\;+\;\)
Hardware-level Control
}

\end{center}

\end{examplebox}

%================================================
% Why Middle Level
%================================================
\booksubheading{C-কে Middle Level Language বলা হয় কেন?}

C ভাষায় একদিকে High Level Programming-এর বিভিন্ন সুবিধা রয়েছে,
অন্যদিকে Low Level Programming-এর মতো Hardware ও Memory-এর
কাছাকাছি কাজ করার সুযোগও রয়েছে।

\begin{tabularx}{\linewidth}{|Y|Y|}
\hline

\rowcolor{softBlue}
\multicolumn{1}{|c|}{\textbf{High Level বৈশিষ্ট্য}}
&
\multicolumn{1}{c|}{\textbf{Low Level বৈশিষ্ট্য}}
\\
\hline

Variable ব্যবহার করা যায়
&
Memory Address নিয়ে কাজ করা যায়
\\
\hline

\textenglish{\ttfamily if, for, while} ইত্যাদি ব্যবহার করা যায়
&
Pointer ব্যবহার করা যায়
\\
\hline

Function তৈরি করা যায়
&
Bitwise Operation করা যায়
\\
\hline

Structured Program লেখা যায়
&
Hardware/System Resource-এর কাছাকাছি কাজ করা যায়
\\
\hline

\end{tabularx}

এই দুই ধরনের বৈশিষ্ট্যের সমন্বয়ের কারণেই C-কে প্রচলিত HSC
শ্রেণিবিভাগে Middle Level Language বলা হয়।

%================================================
% Features
%================================================
\booksubheading{Middle Level Language-এর বৈশিষ্ট্য}

\begin{tabularx}{\linewidth}{|>{\bfseries}p{0.28\linewidth}|Y|}
\hline

\rowcolor{softBlue}
\multicolumn{1}{|c|}{\textbf{বৈশিষ্ট্য}}
&
\multicolumn{1}{c|}{\textbf{ব্যাখ্যা}}
\\
\hline

High ও Low Level-এর সমন্বয়
&
দুই ধরনের Programming বৈশিষ্ট্যই পাওয়া যায়।
\\
\hline

তুলনামূলক সহজ
&
Machine বা Assembly Language-এর তুলনায় Program লেখা ও বোঝা সহজ।
\\
\hline

Hardware Control
&
Memory, Bit ও System Resource-এর কাছাকাছি কাজ করা যায়।
\\
\hline

General Programming
&
সাধারণ Problem Solving ও Application তৈরিতে ব্যবহার করা যায়।
\\
\hline

System Programming
&
System Software ও Hardware-ঘনিষ্ঠ Program তৈরিতে ব্যবহারযোগ্য।
\\
\hline

Efficiency
&
Resource-এর ওপর তুলনামূলক বেশি নিয়ন্ত্রণ থাকায় দক্ষ Program
তৈরি করা সম্ভব।
\\
\hline

\end{tabularx}

%================================================
% Uses
%================================================
\booksubheading{Middle Level Language-এর ব্যবহার}

Middle Level বৈশিষ্ট্যসম্পন্ন ভাষা বিভিন্ন ক্ষেত্রে ব্যবহার করা যায়।
যেমন—

\begin{compacttwocolpoints}

\item System Software তৈরিতে।

\item Operating System-এর বিভিন্ন অংশ তৈরিতে।

\item Embedded System Programming-এ।

\item Hardware-ঘনিষ্ঠ কাজের ক্ষেত্রে।

\item General Purpose Application তৈরিতে।

\item Performance-sensitive Program তৈরিতে।

\end{compacttwocolpoints}

%================================================
% Advantages
%================================================
\booksubheading{Middle Level Language-এর সুবিধা}

\begin{compacttwocolpoints}

\item High Level Language-এর মতো তুলনামূলক সহজে Program লেখা যায়।

\item Low Level-এর মতো Hardware-এর কাছাকাছি কাজ করা যায়।

\item Memory ও System Resource-এর ওপর তুলনামূলক বেশি নিয়ন্ত্রণ পাওয়া যায়।

\item General Programming ও System Programming—উভয় ক্ষেত্রেই
ব্যবহার করা যায়।

\item Function ও Structured Programming ব্যবহার করে বড় Program
সুশৃঙ্খলভাবে লেখা যায়।

\item দক্ষ ও দ্রুত Program তৈরিতে সুবিধা হয়।

\end{compacttwocolpoints}

%================================================
% Limitations
%================================================
\booksubheading{Middle Level Language-এর সীমাবদ্ধতা}

\begin{itemize}

\item সাধারণ High Level Language-এর তুলনায় শেখার কিছু অংশ কঠিন হতে পারে।

\item Pointer ও Memory নিয়ে কাজ করার সময় ভুল হলে Program-এ
জটিল Error তৈরি হতে পারে।

\item Hardware-ঘনিষ্ঠ Programming-এর জন্য Computer Architecture
সম্পর্কে কিছু ধারণা প্রয়োজন হতে পারে।

\item Low-level বৈশিষ্ট্য বেশি ব্যবহার করলে Program-এর
Portability কমে যেতে পারে।

\end{itemize}

%================================================
% Comparison
%================================================
\booksubheading{Low Level, Middle Level ও High Level Language-এর পার্থক্য}

\begin{tabularx}{\linewidth}{|>{\bfseries}p{0.18\linewidth}|Y|Y|Y|}
\hline

\rowcolor{softBlue}
\multicolumn{1}{|c|}{\textbf{বিষয়}}
&
\multicolumn{1}{c|}{\textbf{Low Level}}
&
\multicolumn{1}{c|}{\textbf{Middle Level}}
&
\multicolumn{1}{c|}{\textbf{High Level}}
\\
\hline

Hardware-এর ঘনিষ্ঠতা
&
খুব বেশি
&
প্রয়োজন অনুযায়ী বেশি
&
তুলনামূলক কম
\\
\hline

Program লেখা
&
কঠিন
&
তুলনামূলক সহজ
&
সহজ
\\
\hline

Hardware Control
&
বেশি
&
ভালো
&
তুলনামূলক কম
\\
\hline

General Programming
&
কঠিন
&
উপযোগী
&
খুব উপযোগী
\\
\hline

উদাহরণ
&
Machine, Assembly
&
C, C++
&
Python, Java
\\
\hline

\end{tabularx}

%================================================
% Important Concept
%================================================
\booksubheading{Middle Level Language ও Generation}

Middle Level Language কোনো Programming Language Generation-এর
নাম নয়।

অর্থাৎ—

\begin{center}

\textbf{
Low / Middle / High Level
\(\;\neq\;\)
1GL / 2GL / 3GL / 4GL
}

\end{center}

প্রথমটি মূলত ভাষার বৈশিষ্ট্য ও Hardware-এর সঙ্গে সম্পর্ক বোঝায়,
আর দ্বিতীয়টি Programming Language-এর বিকাশের Generation বোঝায়।

%================================================
% Final Reminder
%================================================
\begin{rememberbox}

\textbf{\color{deepOrange}মনে রাখবে}

\begin{itemize}

\item Middle Level Language-এ
\textbf{High Level ও Low Level—উভয় বৈশিষ্ট্য} পাওয়া যায়।

\item HSC পর্যায়ে \textbf{C Language} সবচেয়ে গুরুত্বপূর্ণ উদাহরণ।

\item C, C++ ইত্যাদি Middle Level বৈশিষ্ট্যসম্পন্ন ভাষার
উদাহরণ হিসেবে উল্লেখ করা হয়।

\item Pointer, Memory Address ও Bitwise Operation হলো
Low Level বৈশিষ্ট্যের উদাহরণ।

\item Variable, Loop ও Function হলো High Level বৈশিষ্ট্যের উদাহরণ।

\item Middle Level Language কোনো
\textbf{Programming Language Generation নয়}।

\end{itemize}

\end{rememberbox}

\small

%================================================
% 5.6 High Level Language
%================================================
